\documentclass[11pt]{article}
\usepackage[utf8]{inputenc}
\usepackage[T2A]{fontenc}
\usepackage[english,russian]{babel}
\usepackage{amsmath,amssymb,bm}
\usepackage[margin=25mm]{geometry}
\usepackage{graphicx}
\usepackage{hyperref}
\title{Counterrotating Motion of Charged Particles in a Two-Frequency Electrodynamic Trap: Dimensionless Theory and Physical Mechanism}
\author{Counterrotating Motion research}
\date{November 2026}
\begin{document}
\maketitle
\begin{abstract}
Рассматривается долговременное встречное вращение двух частиц в идеализированном поперечном восьмигармоническом поле двухчастотной 24-электродной ловушки. Мы восстанавливаем безразмерные уравнения V20B, показываем пределы перенормировки через Reynolds и формулируем контролируемое трёхпараметрическое сечение. Анализ M3 и ведущего отклика объясняет знаки моментов конкуренцией бегущих мод $k=+2$ и $k=-4$, а разные радиусы — балансом RF и DC сил при различных $q/m$. Численные проверки подтверждают формулы в узкой окрестности опубликованного состояния. Они не доказывают существование притягивающего тора и не подтверждают конечную физическую геометрию ловушки.
\end{abstract}
\section{Постановка}
Основной пример --- refine12\_sobol17 при $A=0.0349241$, $B=0.156317$, $\nu=3.21968$, $\beta=0.0153886$, $m_1/m_2=0.109638$, $C_{12}=9.7473\cdot10^{-6}$, $\lambda_1=0.244523$, $\lambda_2=0.177366$. Восемь MFS индексов $k=(-46,-22,2,26,-28,-4,20,44)$ принадлежат двум RF семействам.
\section{Безразмерные уравнения}
Полагая $\tau=\Omega_1t$, $\boldsymbol\xi_i=\mathbf r_i/R$, $s=m_1/m_2$, $Q=q_2/q_1$, получаем $\alpha_1=1$, $\alpha_2=sQ$:
\begin{align}
\boldsymbol\xi_i''={}&\alpha_i\mathbf E_{\rm RF}
+\alpha_i\beta(\xi_x,\xi_y,-2\xi_z)
+\chi_iC_{12}\frac{\boldsymbol\xi_i-\boldsymbol\xi_j}{|\boldsymbol\xi_i-\boldsymbol\xi_j|^3}\\
&-\lambda_i[1+0.15(Re_{*,i}|\boldsymbol\xi_i'|)^{0.687}]\boldsymbol\xi_i'+g_z\mathbf e_z,
\end{align}
где $\chi_1=1$, $\chi_2=s$, $\lambda_i=6\pi\eta a_i/(C_{c,i}m_i\Omega_1)$ и $Re_{*,i}=2\rho_ga_iR\Omega_1/\eta$. Массовый множитель второй частицы обязателен в Coulomb ускорении.
\section{Reynolds координаты и параметры}
Для $\boldsymbol\zeta_i=Re_{*,i}\boldsymbol\xi_i$ аргумент drag становится $|\boldsymbol\zeta_i'|$, но RF гармоника степени $n$ переносит коэффициент $Re_{*,i}^{2-n}$, а Coulomb зависит от $\boldsymbol\zeta_1/Re_{*,1}-\boldsymbol\zeta_2/Re_{*,2}$. Перемена координат не устраняет индивидуальные физические группы. Точное универсальное сокращение к трём параметрам несовместимо с четырьмя независимыми управлениями $A,B,\nu,\beta$. Здесь используется ограниченное сечение $(A,B,\beta)$ при фиксированных остальных группах и координаты $\Pi_T=B^2/A^2$, $\Pi_D=A^2/\lambda_1$, $\Pi_R=\beta/A^2$.
\section{RF моменты}
Для бегущей моды положим $n=|k|$, $\sigma=\operatorname{sgn}k$, $K=\nu_j^2+\alpha\beta$, $D=K^2+(\lambda\nu_j)^2$, $P=\alpha^2n^2|a_jc_k|^2r^{2n-2}$. В линейном ответе при $\omega=0$ полный RF момент равен $\sigma nP\lambda\nu_j/D$. Часть, переданная медленной орбите после учёта момента drag на микродвижении, равна $\sigma(n-1)P\lambda\nu_j/D$. Поэтому нуль полного RF момента и нуль момента медленного движения не совпадают.
\section{Радиальная локализация}
Для ведущих мод $+2,-4$ средняя RF радиальная сила:
\[
\bar F_r=-4\alpha^2A^2|c_2|^2\frac{1+\alpha\beta}{D_1}r
-48\alpha^2B^2|c_{-4}|^2\frac{\nu^2+\alpha\beta}{D_2}r^5.
\]
Она уравновешивает наружную DC силу $\alpha\beta r$. Поскольку RF вклад квадратичен по $\alpha$, а DC линейный, частица с меньшим отношением заряда к массе локализуется дальше от оси. Ведущая модель даёт $r/R=0.55745,0.75167$ против полного RF результата $0.55572,0.75152$.
\section{M3 проверка и моментный баланс}
На M3 корне сетки $192\times16\times16$ пересчёт восьми мод даёт для частицы 1 моменты RF семейств $+7.5171\cdot10^{-4}$ и $-3.6499\cdot10^{-4}$; для частицы 2 $+1.6793\cdot10^{-4}$ и $-2.6012\cdot10^{-4}$. Суммарные RF моменты компенсируются нелинейным drag. Допуски fast HB residual составили $4.1\cdot10^{-9}$ и $1.6\cdot10^{-10}$.
\section{Coulomb и орбитальная киральность}
Опубликованный Coulomb sweep сохраняет встречные знаки при $C=0$ и при номинальной связи. Следовательно, Coulomb не является необходимым источником самого встречного вращения. Знак $s_i=\operatorname{sgn}\omega_i$ можно использовать как орбитальную киральность; это не квантовый спин. Синхронизация фаз и кристаллический порядок не установлены.
\section{Устойчивость и границы вывода}
V20B сохраняет встречное движение на 240000 RF периодах. M3 slow flow имеет локально отрицательные собственные значения, но два условных показателя полной механической системы близки к нулю. Долгое удержание и устойчивость усреднённого корня не доказывают притяжение к инвариантному тору. Кроме того, RF модель использует бесконечные стержни и идеальное DC поле.
\section{Заключение}
Механизм направленного момента и различия радиусов количественно согласуется с конкуренцией бегущих RF гармоник и с балансом RF/DC сил. Трёхпараметрическое представление оправдано лишь для явно фиксированного сечения. Следующий динамический вопрос --- независимые отключения RF семейств и их влияние на устойчивый M3 корень.
\end{document}
